\chapter{Packet Analysis and Performance Metrics}
\label{ch:analysis}

\begin{defbox}[In this chapter you will learn]
\begin{itemize}
 \item the standard metrics---delivery error rate (DER), throughput, per-node and per-base-station statistics---and exactly how the demos and the analysis module define them;
 \item the \emph{two} ways Amber classifies collisions (real-time and post-hoc) and why they can disagree;
 \item a verified defect in the post-hoc path that makes almost every packet look collided;
 \item practical guidance on which numbers to trust and how to report them.
\end{itemize}
\end{defbox}

\section{Metrics}
\label{sec:metrics}

A simulator is only as useful as the questions its output can answer. For Ambient IoT networks the natural metrics are the following.

\begin{defbox}[Delivery and throughput]
Let $N_{\mathrm{sent}}$ be the number of packets nodes handed to the channel, $N_{\mathrm{dec}}$ the number decoded without collision loss at the base station, and $T$ the simulated time.
\[
 \mathrm{PDR}=\frac{N_{\mathrm{dec}}}{N_{\mathrm{sent}}},\qquad
 \mathrm{DER}=1-\mathrm{PDR},\qquad
 \text{throughput}=\frac{N_{\mathrm{dec}}}{T}\ [\text{packets/s}].
\]
\end{defbox}
Other useful quantities: per-node success counts (fairness), the number of registered nodes versus time (discovery), collisions per frame, and energy-related ones such as the number of FSM cycles per charge cycle or the fraction of time a node is active.

\subsection{How the demos compute them}

Each demo prints
\[
 \texttt{total\_sent}=\sum_{\text{nodes}}\codew{packets_sent},\qquad
 \texttt{decoded}=\bigl|\{(\text{node},\text{window start}):\ \neg\text{collided}\}\bigr|,
\]
\[
 \texttt{DER}=1-\texttt{decoded}/\texttt{total\_sent},\qquad
 \texttt{throughput}=\texttt{decoded}/(\texttt{SIM\_TIME\_MS}/1000).
\]
Note that the numerator counts \emph{unique} (node, window) pairs, and that ``sent'' counts all handed-over packets including those the base station later drops for being below its sensitivity (the denominator is not reduced by them). For the unicast demo in one run: $125$ sent, $125$ decoded, DER $=0\%$, $6.2$ packets/s.

\subsection{Name clash: ``DER''}

\begin{reality}[Two different quantities share the name ``DER'']
In the demos DER is an \emph{error} rate: $1-\text{decoded}/\text{sent}$ (smaller is better; the plots label it ``Delivery Error Rate''). In \codew{packet_analysis.analyze_packets}, \code{BSStats.der} is $\text{received}/\text{transmitted}$, a \emph{success} ratio (larger is better), and ``transmitted'' there means ``rows in the base station's RX log'', i.e.\ packets \emph{logged} by the base station, not packets \emph{sent} by nodes. In our test run the same simulation gives demo-style DER $\approx2.4\%$ (121 decoded of 124 sent) but analysis-module \code{der}$=0.976$. Always state which you mean.
\end{reality}

\section{The \texttt{packet\_analysis} module}
\label{sec:pamodule}

Besides \code{apply_sic} and the noise helper used by the engine, \file{packet_analysis.py} provides a post-processing pipeline.

\begin{center}
\small
\renewcommand{\arraystretch}{1.12}
\begin{tabularx}{\linewidth}{@{}>{\ttfamily}p{4.0cm}X@{}}
\toprule
\normalfont\bfseries Function / class & \normalfont\bfseries Role\\
\midrule
load\_packet\_log(path) & read a tab-separated BS RX log; requires \codew{start_ms, end_ms, node_id, rssi_dbm, subcarrier_shift}\\
\codew{detect_collisions(df,...)} & \emph{post-hoc} collision detection by time overlap, then SIC per overlap\\
analyze\_packets(\dots) & per-BS and per-node statistics; writes optional text logs\\
BSStats & \code{transmitted}, \code{received}, \code{errors}, \codew{collisions_total}, \code{der}\\
NodeStats & \code{sent}, \codew{received_at_bs}, \codew{total_errors}, \code{heard_by_bs}\\
CollisionEvent & kind $\in\{\text{total},\text{n\_way},\text{capture},\text{sic\_success}\}$, nodes, powers, SINRs, winners\\
\bottomrule
\end{tabularx}
\end{center}

\paragraph{How \code{analyze_packets} chooses its path.} If the log has a \code{collided} column (which the base station writes) the function \emph{trusts it}: received $=$ rows with \codew{collided==False}, errors $=$ rows with \codew{collided==True}. Only if the column is missing does it call \codew{detect_collisions} and re-derive the outcomes.

\paragraph{Post-hoc detection.} For each subcarrier it builds start/end events for each packet interval, sweeps in time order, and whenever a packet starts while others are active, it forms a collision set and applies SIC (default $\gamma_{\mathrm{req}}=10$~dB here, versus $3$~dB in the engine) to label packets \code{decoded}, \code{captured}, \codew{sic_recovered} or \codew{collision_loss}. Collision events are classified \code{total} (two packets, none decoded), \code{n_way} ($>2$, none decoded), \code{capture} (one decoded) and \code{sic_success} (two or more).

\paragraph{Side effects to know about.} By default \codew{analyze_packets} \emph{appends} a text report to \codew{terminal_out.txt} and collision lines to \codew{collision_out.txt} in the current directory on every call; pass \codew{write_terminal_log=None, write_collision_log=None} to disable. The demos' own CSV files (\codew{rx_packets.csv}) lack \code{start_ms}/\code{end_ms}, so \codew{load_packet_log} rejects them (``Missing required columns''): the analysis module expects logs written by the engine's \code{save_logs}/\code{stream_log}.

\section{Real-time versus post-hoc classification}
\label{sec:rtposthoc}

We ran the same $\SI{20}{s}$ broadcast simulation (SIC off) and analysed its engine-written log in both ways.

\begin{center}
\small
\renewcommand{\arraystretch}{1.12}
\begin{tabular}{@{}lrrr@{}}
\toprule
Path & Received (decoded) & Errors & Collision events\\
\midrule
Real-time (\code{collided} column) & 121 & 3 & --\\
Post-hoc, as logged & \textbf{1} & \textbf{123} & 122\\
Post-hoc, zero-length intervals nudged by $\SI{1}{\micro s}$ & 120 & 4 & 4\\
\bottomrule
\end{tabular}
\end{center}
The real-time classification and the (repaired) post-hoc classification agree to within one packet; the unrepaired post-hoc path reports that $99\%$ of packets were lost.

\begin{reality}[A latent defect in the post-hoc sweep]
In the engine's log, a packet's \code{start_ms} is the start of the RX \emph{window} and \code{end_ms} is its \emph{arrival} time. A node that happens to be ready when the window opens transmits in the very first millisecond, so \codew{start_ms == end_ms}: $100$ of the $124$ packets in our run have zero-length intervals. The sweep sorts events by \codew{(time, is_start)}, i.e.\ \emph{end} events before \emph{start} events at the same time. For a zero-length packet the \emph{end} is therefore processed first (removing nothing, marking the packet \code{decoded}) and the \emph{start} afterwards, which inserts the packet into the \code{active} set \emph{permanently}. Every later packet then ``collides'' with this ghost, producing the cascade of $n$-way collisions ($122$ events in our run). The one-line remedy (end events after start events for equal times, or a minimum duration) restores agreement. It does not affect the demos, which count via the engine's real-time flags, but it makes \codew{analyze_packets} unreliable on logs lacking the \code{collided} column. Remember also that, for engine logs, the overlap intervals are not airtime: they are ``window start to arrival'', so even a repaired post-hoc detector approximates collision by an interval-overlap rule different from the engine's $\SI{5}{ms}$ rule.
\end{reality}

\section{Practical guidance}

\begin{enumerate}
 \item \textbf{Prefer the engine's counters} (\code{total_rx}, \codew{total_collided}, \codew{rx_count_by_node}) and \code{collided} flags for headline results; they are cheap, constant-memory and come from the classification actually used by the simulation.
 \item \textbf{Report what you divide by.} ``Sent'' includes packets later dropped below sensitivity; if you want the PDR of \emph{audible} packets, count rows in \code{rx_packets} instead.
 \item \textbf{Report distributions.} Replicate and show box plots or confidence intervals (Chapter~\ref{ch:des}); a single unseeded run is one sample.
 \item \textbf{Separate energy from radio losses.} A silent node is not a lost packet. Count \emph{opportunities} (polls to a node that was in \emph{wait\_slot}) if you want a radio-only PDR.
 \item \textbf{Mind the logging overhead.} For long runs use \codew{keep_logs=False} and \code{stream_log}.
\end{enumerate}

\begin{codemap}
\small
\begin{tabularx}{\linewidth}{@{}>{\ttfamily}p{5.0cm}X@{}}
\toprule
\normalfont\bfseries Concept & \normalfont\bfseries Code\\
\midrule
Demo metrics & \codew{run_simulation()} in each demo; \codew{save_packets}\\
Always-on counters & \codew{BSBehavior.total_rx}, \codew{total_collided}, \code{total_tx}, \codew{rx_count_by_node}\\
Offline statistics & \codew{analyze_packets} (\file{packet_analysis.py})\\
Offline collisions & \codew{detect_collisions}, \codew{_analyze_collision_outcome}\\
\bottomrule
\end{tabularx}
\end{codemap}

\section{Chapter summary}
\begin{itemize}
 \item DER (error) $=1-\text{decoded}/\text{sent}$ in the demos; \code{BSStats.der} is a success ratio over logged receptions. Do not mix them.
 \item Prefer the engine's real-time \code{collided} flags; the post-hoc detector is unreliable on engine logs because of zero-length intervals (we measured $1$ vs $121$ decoded).
 \item \codew{analyze_packets} writes text files by default and does not read the demos' CSV files.
 \item Report distributions, and separate energy-induced silence from channel-induced loss.
\end{itemize}

\section*{Exercises}
\begin{exercise}
A run reports $\texttt{total\_sent}=130$, $\texttt{decoded}=117$ over $20.001$~s. Compute DER, PDR and throughput. If $6$ of the missing $13$ packets were dropped below sensitivity, what is the PDR of audible packets?
\end{exercise}
\begin{exercise}
Explain why \codew{start_ms == end_ms} for a packet that arrives in the first millisecond of its window, and sketch the event order of the sweep for two such packets in consecutive windows.
\end{exercise}
\begin{exercise}
Design a metric that separates ``lost to energy'' from ``lost to collisions'' using only data the base station can log. Which additional information from the simulator (not available in a real network) would you use as ground truth?
\end{exercise}
